\section{Why Schedule the Complete Pipeline?}
\label{sec:legacy-motivation}
The original real FFT suspends a conventional ordered butterfly traversal.
With $B$ butterflies and requested horizon $H$, it executes
$q=\ceil{B/H}$ butterflies per sample call and completes after
$L=\ceil{B/q}\le H$ calls. The original modules immediately restart completed
transforms, so their steady-state frame interval is $L$, which need not equal
$H$. For $N=4096,H=2048$, $L=1878$: the frame rate is about 9.1\% faster
than requested, and temporal smoothing based on the nominal hop decays too
quickly. Appendix~\ref{sec:scheduling} retains the completion proof, cadence
table, timestamp diagram, and smoothing derivation.

Budgeting butterflies also leaves substantial boundary work. The original
path packs and permutes a frame in bulk, reconstructs all real-spectrum bins
on the completing call, and processes magnitude/display arrays in separate
passes; some of those paths allocate. For $H=\rho N$ with fixed $\rho>0$,
the butterfly quota grows as $O(\log N)$, while preparation and reconstruction
still impose $O(N)$ boundary passes. Appendix~\ref{sec:cost} gives the cost
model and ownership analysis. These two limitations motivate an exact frame
clock and a work sequence covering the complete analysis, not just its FFT.
